\section{Introduction}
\label{sec:intro}

Embodied AI and robot-learning systems are typically evaluated through benchmark pipelines rather than through a model in isolation. A single run may depend on the policy checkpoint, simulator, task suite, reset state, controller semantics, observation preprocessing, dependency versions, seeds, and success definition. Prior work shows that reinforcement-learning results are sensitive to evaluation and experimental choices \cite{henderson2018drlmatters}; in embodied settings, those choices are embedded in larger evaluation pipelines. As a result, a changed success rate or episode outcome is often ambiguous: it may indicate a likely true policy regression, a setup-sensitive evaluation change, or evidence too weak to support attribution.

This ambiguity creates a practical diagnosis problem. Aggregate metrics can show that a benchmark result changed, but not why. Blind reruns can test repeatability, but do not identify whether the cause is a controller-interface mismatch, a changed reset state, an observation adapter, dependency drift, or a policy regression. Manifest differences identify changed fields, but a diff alone cannot establish causality. Interpreting embodied-AI evaluations therefore requires both structured provenance for candidate factors and behavioral or failure-recovery evidence to test them.

\sys addresses this problem by treating evaluation deviation diagnosis as a factor-directed evidence task. For baseline, current, and replay runs, it records structured manifests, summaries, episode outcomes, logs, and failure records when execution fails. It detects aggregate, paired-outcome, and failed-run deviations; maps manifest and failure differences to factor families such as action/controller interface, observation preprocessing, reset state, protocol/metric, dependency/runtime, checkpoint/config compatibility, and data format; and schedules targeted differential replays. These factors are hypotheses, not diagnoses. \sys attributes a setup-sensitive deviation only when replay or recovery evidence supports the factor, reports likely true regression when external-factor evidence does not explain a thresholded deviation, and returns unknown when the evidence is weak.

We ground the factor space through an empirical analysis of real project records. Across 15 embodied AI and robot-learning projects, keyword mining produced 2,714 candidate issues and pull requests, a relevance screen assisted by a large language model (LLM) retained 641, and manual classification retained 473 evidence records for the final taxonomy. We then evaluate \sys on a LeRobot--LIBERO artifact corpus covering completed rollouts, failed executions, calibration cases, ablations, robustness scenarios, and true-regression probes. \sys attributes all detected paper-only and robustness cases in our corpus to the expected factor, makes no false setup attributions on 12 negative calibration cases, classifies 66/67 main-status cases, and reduces diagnosis workload by 69.3\% compared with a blind rerun-$k$ baseline.

This paper makes the following contributions:
\begin{itemize}
  \item We frame embodied AI evaluation deviation diagnosis as an evidence-aware benchmark interpretation problem that complements aggregate metrics with evidence.
  \item We present a taxonomy of deviation symptoms and engineering factors from 473 GitHub issues and pull requests across 15 embodied AI and robot-learning projects.
  \item We introduce \sys, a manifest- and replay-based diagnosis technique for deviation detection, factor-directed replay, attribution, and abstention.
  \item We evaluate \sys on LeRobot+LIBERO artifacts spanning completed-rollout deviations, failed-run cases, negative calibration, ablations, robustness cases, true-regression probes, and replay-cost measurements.
\end{itemize}
